\subsection{Marked period components and absolute unfolding}

Fix a rational Hermitian space \(W\) over \(E/F\), \(F=E^+\), its
polarizing trace form and CM type. Take the identity auxiliary cover
and its principal odd level \(M\), with
\((M)=\mathfrak m\bar{\mathfrak m}\) and
\(\mathfrak m+\bar{\mathfrak m}=O_E\).
In the period description a component has a self-dual row lattice
\(L\subset W\), the polarized half-level mark and the connected flag
domain of its prescribed signature. A marking on one half of \(L/ML\)
can equally be specified by its dual basis on the opposite half,
using the fixed primitive \(M\)-th root for the Weil pairing.

\begin{lemma}\label{munf:marks}
Fix a reference row \(L_*\) with a basis of \(L_*/\mathfrak mL_*\).
Every such marked period lattice in \(W\) can be carried by a rational
isometry to a row equal to \(L_*\) at the primes of \(M\), with the
specified basis on the \(\mathfrak m\)-half. Write \(\Delta\) for the
rational unitary elements integral there and acting identically on this
marked half of \(L_*\). Two transported rows lie in the same
\(\Delta\)-orbit exactly when their polarized marked period lattices
are isometric. The stabilizer of each row is its principal-level
arithmetic group.
\end{lemma}
\begin{proof}
Every place \(v\) of \(F\) above the level splits in \(E\), since a
prime fixed by conjugation could not divide the product of the
coprime conjugate halves. The polarization pairs the two summands of
\(L/ML\) perfectly: the CRT idempotents are interchanged by the
adjoint involution, so each summand is isotropic and is the dual of the
other. This proves the equivalence of the two ways of specifying
the marking, including at a prime power. Locally \(U(W)(F_v)\) is
\(\mathrm{GL}_r(F_v)\) on the first factor, acting by the inverse
adjoint on the second. Lift the indicated residue bases in both
first-factor lattices to bases over the DVR. The matrix carrying one
basis to the other carries the lattices onto one another. Its inverse
adjoint carries the second lattices and the dual marks, by perfection.

The resulting conditions on the tuple of local isometries form a
nonempty open set. Rational isometries meet it. Indeed the Cayley
map \(X\mapsto(1-X)(1+X)^{-1}\) identifies an open of the
\(F\)-vector space of skew-adjoint endomorphisms with the unitary
open where \(1+g\) is invertible. Every selected local open meets
this chart: multiply its matrix by a norm-one scalar arbitrarily
close to one, avoiding the finitely many prohibited eigenvalues.
At these split places such scalars are \((t,t^{-1})\).
Weak approximation on the vector space, in a fixed \(F\)-basis,
now gives a single rational isometry. The transported lattices agree
with \(L_*\) at the whole level. Globally, this is equivalent to the
existence of a common denominator prime to the level in the two
lattice inclusions: localize the finite torsion modules
\((L+L_*)/L\) and \((L+L_*)/L_*\), then choose an annihilator
congruent to one at those primes by CRT.

An element of \(\Delta\) carrying one transported lattice to another
preserves the mark on the first half, hence on the dual half.
Conversely a marked isometry between the lattices extends uniquely
by \(E\)-linearity to \(W\); its local action belongs to \(\Delta\).
When the two rows agree this gives precisely the integral isometries
acting identically modulo \(M\). To see the component statement in
period coordinates, a complex framed polarized torus has a period
lattice on its tangent space. An isomorphism lifts uniquely to its
complex-linear derivative on the universal covers, so it induces
the preceding integral isometry. Conversely transporting the
complex structure and flag by an isometry gives isomorphic marked
complex tori. For a fixed marked lattice the positive-plane domains
are matrix balls with connected compact flag fibres; their product
is connected. In a family of complex abelian varieties the integral
homology, its trace pairing and the finite marking are locally
constant in a topological trivialization. This gives exactly the
isometry classes as component indices in the marked period model.
\end{proof}
\src{src/Development/AbelianVarietyTheory/Proofs/SiegelThetaTransformation/ExistsMoveMarked.lean}{133}{Rational transport of the marked basis},
\src{src/Development/AbelianVarietyTheory/Proofs/SiegelThetaTransformation/ComponentEqOfMatchedLatticesMarkings.lean}{74}{identification of marked indices} and
\src{src/Development/AbelianVarietyTheory/Proofs/SiegelThetaTransformation/ExistsComponentsLatticesMovedToLevelFamily.lean}{162}{distinct orbits after simultaneous transport}.

\begin{proposition}[Unfolding the finite span]\label{munf:span}
On the countable set of admissible rows let \(C\) be the equivariant
matrix of Lemma~\ref{bfam:gram} and \(K\) the kernel of
Proposition~\ref{arrep:reproduce}. Suppose the marked periodized
kernel is the finite sum of paired theta series for the indicated
column genus, with a single nonzero normalization. Then every
bounded family of holomorphic forms of the prescribed weight and
unitary multiplier on these row orbits lies in the span of those
theta series. The conclusion also holds simultaneously on any finite
collection of distinct row orbits.
\end{proposition}
\begin{proof}
Write the input as \(F\) and put \(V=C^{-1}F\). Use the norm given
by the supremum over all rows and orthoframes. The Neumann series
converges in this norm and preserves the weight and multiplier,
since \(C\) commutes with the action. It converges locally uniformly
on the full domain by triangular equivariance and Gram--Schmidt,
so \(V\) is holomorphic there. Its absolute value is bounded on all
rows of the slab by one constant \(B\), by the compactness of the
triangular factors in the proof of Proposition~\ref{arrep:reproduce}.
For a fixed output \(a=(L,x)\) put
\[
          f_a(L',y)=C_{LL'}K(x,y)V_{L'}(y).
\]
On rows times the slab we have
\[
 \sum_{L'}\int_{S_\delta}|f_a(L',y)|\,dy
       \le B\sum_{L'}|C_{LL'}|\int_{S_\delta}|K(x,y)|\,dy<\infty .
\]
Thus we may unfold under the countable diagonal action on this
product with its counting and Lebesgue measures. The action
preserves both measures.

Here is an explicit reason a measurable fundamental set exists.
Choose one \(L'\) in each row orbit. Its stabilizer is discrete in
the real unitary group: on an integral basis of the underlying
rank-\(2fr\) lattice it has integral matrix entries. It acts freely
on the invertible frames and properly, since the isometries carrying
one compact set of frames into another form a subset with compact
closure (write \(g=y'y^{-1}\)). Cover the slab by countably many
open patches disjoint from their nonidentity stabilizer translates.
Keep in each patch the complement of the translates of all earlier
patches. Their union meets every stabilizer orbit once and is Borel.
The union on the representative rows is a diagonal fundamental set
\(\mathcal D_0\). Absolute integrability and Tonelli now give
\[
 \sum_{L'}\int_{S_\delta}f_a(L',y)\,dy
       =\int_{\mathcal D_0}\sum_{\gamma\in\Delta}
                           f_a(\gamma L',\gamma y)\,d(L',y).
\]
The left side is \(\rho_\delta(CV)_L(x)=\rho_\delta F_L(x)\).
On the right the multiplier law for \(V\) and the given periodized
identity express the inner sum as \(\sum_{j=1}^s q_j(a)h_j(L',y)\),
where the \(q_j\) form a basis selected from the finitely many theta
series. The functions \(h_j\) are integrable even if a single
uncombined term in the original genus expression has no known
integral. Indeed choose \(s\) evaluations \(a_i\) for which
\((q_j(a_i))\) is invertible: all evaluation vectors span
\(\mathbf C^s\) by independence of the functions. The sums at these
\(s\) points are integrable by the preceding absolute estimate, and
matrix inversion expresses each \(h_j\) as their linear combination.
The zero-dimensional span gives zero directly. Divide by the
positive \(\rho_\delta\) to obtain the asserted finite expansion.
For a collection of orbits extend the families by the equivariant
law from each representative and by zero on the complement. The
stabilizer law makes the extension well-defined and the unitary
multiplier preserves its bound.
\end{proof}
\src{src/Development/AbelianVarietyTheory/Proofs/SiegelThetaTransformation/ExistsFundamentalDomain.lean}{138}{A fundamental set on lattice--frame pairs} and
\src{src/Development/AbelianVarietyTheory/Proofs/SiegelThetaTransformation/HermitianThetaSeriesSpanBoundedForms.lean}{22}{finite rank after absolute unfolding}.

For the rows of Lemma~\ref{munf:marks} in their common rational
Hermitian space, use the graph sum of Proposition~\ref{dg:identity}.
The passage to algebraic cusp sections then uses \eqref{bfam:reading}
on that space with its marked multiplier.

